\documentclass[11pt,a4paper]{article}
\usepackage[T1]{fontenc}
\usepackage{amsmath,amssymb,amsthm}
\usepackage[margin=27mm]{geometry}
\usepackage[hidelinks]{hyperref}
\newtheorem{proposition}{Proposition}
\newcommand{\Z}{\mathbb Z}
\newcommand{\Q}{\mathbb Q}
\newcommand{\ph}{\varphi}
\newcommand{\Tr}{\operatorname{Tr}}
\newcommand{\Nm}{\operatorname{N}}
\title{Twist Numbers:\\Exact Integer Arithmetic and the J Step}
\author{A. M. Thorn\\\small\href{https://orcid.org/0009-0008-5463-278X}{ORCID 0009-0008-5463-278X}}
\date{18 November 2026}
\begin{document}
\maketitle
\begin{abstract}
Twist numbers are the elements of the classical cyclotomic ring
$\Z[j]$, where $j=e^{2\pi i/5}$. We give its unique four-integer
coordinates, product rule and real subring $\Z[\ph]$.
Multiplication by $J=1+j^2$ and its inverse each admit a schedule of
four integer additions or subtractions. The identities $\ph J=j$
and $(J-1)^3=j$ connect the golden unit, the root of unity and the
step operator. Two complex embeddings show contraction and expansion
with determinant one. These explicit coordinates provide an algebraic
foundation for TWIST-J; they do not by themselves specify a physical
update law or a spacetime interpretation.
\end{abstract}

\section{The ring and its coordinates}\label{sec:coordinates}
Fix the embedding $j=e^{2\pi i/5}$ and put
\[
 \Phi_5(x)=1+x+x^2+x^3+x^4,\qquad
 R=\Z[j]\simeq\Z[x]/(\Phi_5(x)),\qquad K=\Q(j).
\]
The polynomial $\Phi_5(x+1)=x^4+5x^3+10x^2+10x+5$ is Eisenstein
at $5$, so $\Phi_5$ is irreducible over $\Q$. Division by this monic
polynomial gives the unique coordinates
\begin{equation}\label{eq:coordinates}
 \alpha=a+bj+cj^2+dj^3,\qquad (a,b,c,d)\in\Z^4.
\end{equation}
Thus four-dimensional means rank four over $\Z$, or degree four over
$\Q$ after passing to $K$. It specifies neither a spacetime dimension
nor a bound on the size of the four integers.

Addition is coordinatewise. For multiplication, form the ordinary
polynomial product $\sum_{k=0}^6u_kj^k$ and use
$j^4=-1-j-j^2-j^3$, $j^5=1$, $j^6=j$. The result is
\begin{equation}\label{eq:product}
 (u_0-u_4+u_5)+(u_1-u_4+u_6)j+(u_2-u_4)j^2+(u_3-u_4)j^3.
\end{equation}
This supplies exact, commutative and associative arithmetic with integer
coefficients. The name twist numbers denotes this coordinate presentation
of the familiar cyclotomic ring, not a new number field \cite{washington}.

\section{The golden real subring}\label{sec:real}
Define the algebraic element
\begin{equation}\label{eq:golden}
 \ph=-j^2-j^3.
\end{equation}
In the fixed embedding it is $(1+\sqrt5)/2$. Its sign and numerical
interpretation refer to that embedding; a Galois conjugate need not be
positive.

\begin{proposition}\label{prop:real}
One has $\ph^2=\ph+1$, $\ph^{-1}=\ph-1=j+j^4$, and
$R\cap\mathbb R=\Z[\ph]$ in the fixed embedding. Moreover
\[
 R=\Z[\ph]\oplus j\Z[\ph].
\]
\end{proposition}
\begin{proof}
Expanding $(-j^2-j^3)^2$ and reducing powers gives $\ph^2=\ph+1$;
the other scalar identity follows from the cyclotomic relation.
Complex conjugation sends the coordinates in \eqref{eq:coordinates} to
$(a-b,-b,d-b,c-b)$. Equality with $(a,b,c,d)$ is equivalent to
$b=0$ and $c=d$. The fixed elements are therefore exactly $a-c\ph$.
Finally,
\[
 j^2=(\ph-1)j-1,\qquad j^3=(1-\ph)(1+j).
\]
These identities give the displayed decomposition. It is unique because
$r+js=0$, with real $r,s$, forces $s=0$ and then $r=0$.
\end{proof}

The field tower is $\Q\subset\Q(\sqrt5)\subset K$, with both degrees
equal to two. For $r=a+b\ph$, the quadratic norm is
\[
 \Nm_{\Q(\sqrt5)/\Q}(r)=a^2+ab-b^2.
\]
It must be distinguished from the quartic norm:
$\Nm_{K/\Q}(r)=(a^2+ab-b^2)^2$. In particular $\ph$ has
quadratic norm $-1$ and quartic norm $1$.

\section{The distinguished unit}\label{sec:unit}
Set $J=1+j^2$. Multiplication by the closure relation gives
\begin{equation}\label{eq:unit}
 \ph J=j,\qquad J^{-1}=\ph j^{-1}=-j-j^2,\qquad (J-1)^3=j.
\end{equation}
Consequently $J$ is a unit of $R$ and $\Z[J]=R$: $J$ determines
the original generator as well as belonging to its ring. Also
\begin{equation}\label{eq:polar}
 J=\ph^{-1}e^{2\pi i/5},\qquad J^5=\ph^{-5}.
\end{equation}
The polar statement uses the fixed embedding; the fifth-power identity
is an identity in the ring.

For $r\in\{1,2,3,4\}$ let $\sigma_r(j)=j^r$. The field norm is
$\Nm_{K/\Q}(\alpha)=\prod_{r=1}^4\sigma_r(\alpha)$.
The evaluation matrix $(j^{rk})_{1\leq r\leq4,\,0\leq k\leq3}$
is a Vandermonde matrix at four distinct roots, hence invertible;
it diagonalizes multiplication by $\alpha$ with entries $\sigma_r(\alpha)$.
Thus for $\alpha\in R$ the norm equals the integer determinant of
multiplication by $\alpha$ in the integral basis. It is positive for
nonzero $\alpha$, since the factors occur in complex-conjugate pairs.
An element of $R$
is a unit precisely when this determinant is $1$: necessity follows
from multiplicativity, and sufficiency from the integral inverse
matrix applied to the coordinate vector of $1$.

\section{Matrices and the exact step}\label{sec:step}
Use column vectors in the ordered basis $(1,j,j^2,j^3)$. Multiplication
by $j$ and by $J$ has the matrices
\begin{equation}\label{eq:matrices}
 M_j=\begin{pmatrix}
 0&0&0&-1\\1&0&0&-1\\0&1&0&-1\\0&0&1&-1
 \end{pmatrix},\qquad
 M_J=I+M_j^2=\begin{pmatrix}
 1&0&-1&1\\0&1&-1&0\\1&0&0&0\\0&1&-1&1
 \end{pmatrix}.
\end{equation}
Each column follows by multiplying one basis element and applying
\eqref{eq:product}. In particular $M_j^5=I$.

\begin{proposition}\label{prop:step}
The forward and inverse coordinate rules are
\begin{align}
 (a,b,c,d)&\longmapsto(a-c+d,\ b-c,\ a,\ b-c+d),\label{eq:forward}\\
 (A,B,C,D)&\longmapsto(C,\ C+D-A,\ C+D-A-B,\ D-B).\label{eq:inverse}
\end{align}
Each rule uses four additions or subtractions and no multiplications,
when repeated intermediate values are reused. Furthermore,
\begin{equation}\label{eq:charpoly}
 \det M_J=1,\quad\Tr M_J=3,\quad
 \det(xI-M_J)=x^4-3x^3+4x^2-2x+1.
\end{equation}
\end{proposition}
\begin{proof}
Equation \eqref{eq:forward} is \eqref{eq:matrices} applied to a column
vector. Compute $t=a-c$ and $u=b-c$, then output $(t+d,u,a,u+d)$.
For the inverse, compute $t=D-A$, $b=C+t$, $c=b-B$, $d=D-B$ and
output $(C,b,c,d)$. Substitution verifies both compositions.
The trace is read from the matrix. Since squaring permutes the
nontrivial fifth roots, the eigenvalues are the four numbers
$1+j^r$, and
\[
 \det(xI-M_J)=\prod_{r=1}^4(x-1-j^r)=\Phi_5(x-1).
\]
Expansion gives \eqref{eq:charpoly}, including the determinant.
\end{proof}

For example, the unit element makes the exact round trip
\begin{equation}\label{eq:example}
 (1,0,0,0)\xrightarrow{J}(1,0,1,0)
 \xrightarrow{J^{-1}}(1,0,0,0).
\end{equation}

The operation count concerns integer arithmetic operations; it is not
a constant bit-cost claim. Coefficients can grow under iteration.
Likewise $M_j^2$ has order five, but is not orthogonal for the ordinary
Euclidean dot product on coefficient vectors. Its rotation interpretation
belongs to the embedding described next.

\section{Two complex embeddings}\label{sec:embeddings}
The map
\[
 \alpha\longmapsto\bigl(\sigma_1(\alpha),\sigma_2(\alpha)\bigr)
 \quad\hbox{into }\mathbb C^2\simeq\mathbb R^4
\]
is the Minkowski representation, choosing one embedding from each
conjugate pair. The Vandermonde argument above, together with complex
conjugation, makes its real-linear extension from $K\otimes_{\Q}\mathbb R$
an isomorphism. In these coordinates multiplication by $J$ acts
separately on the two complex planes:
\begin{equation}\label{eq:embeddings}
 \sigma_1(J)=\ph^{-1}e^{2\pi i/5},\qquad
 \sigma_2(J)=\ph e^{-\pi i/5}.
\end{equation}
The second formula follows from
$1+e^{-2\pi i/5}=2\cos(\pi/5)e^{-\pi i/5}$ and
$2\cos(\pi/5)=\ph$. The first is \eqref{eq:polar}.

One real two-plane contracts by $\ph^{-1}$ while the other expands
by $\ph$. The real determinant is therefore
$\ph^{-2}\ph^2=1$. The product of the two complex moduli is $1$;
the quartic field norm is the square of that product. Since no
eigenvalue has modulus one, $M_J$ is hyperbolic. Since
$|\sigma_1(J)|<1$, $J$ has infinite order. None of these statements
asserts an increase of every coefficient norm at each step.

\section{Scope and relation to TWIST-J}\label{sec:scope}
These classical algebraic facts provide explicit coordinates for the
J identities used in TWIST-J \cite{canon}. Toral dynamics is a separate
application \cite{hyperbolic}. Public Canon v97 adopts $J=1+\zeta_5^2$
as a primitive axiom of the physical program; its state architecture,
full update law and decoder are additional declarations. The identities
here do not select those declarations or supply a physical clock,
metric, force or occurrence law. In particular $M_J=I+M_j^2$ is a
linear identity, not a decomposition of the full native update into
physical rest and motion.

\paragraph{Verification.}
The accompanying \texttt{verify\_algebra.py} checks the finite algebraic
identities using integer polynomial reduction and independent determinant
expansion. Checks on basis vectors establish the displayed linear maps;
checks on basis pairs establish the product rule by bilinearity. This is
an editorial check of the written proofs, not a new formal Canon probe
or a verification of a physical interpretation.

\paragraph{Edition.}
This manuscript revises the source edition at \href{https://doi.org/10.5281/zenodo.20559777}{Zenodo 20559777}.

\begin{thebibliography}{9}
\bibitem{washington} L. C. Washington, \emph{Introduction to Cyclotomic
Fields}, second edition, Springer, 1997.
\bibitem{hyperbolic} A. M. Thorn, \emph{The Hyperbolic Automorphism J:
A Golden Ratio Unit in $\Z[\zeta_5]$}, 2026.
\href{https://doi.org/10.5281/zenodo.18523754}{doi:10.5281/zenodo.18523754}.
\bibitem{canon} A. M. Thorn, \emph{TWIST-J Public Canon v97},
\href{https://github.com/mathorn1973/twist-j/tree/canon-v97}{tag canon-v97}.
Authority is declared in
\href{https://github.com/mathorn1973/twist-j/blob/canon-v97/STATUS.md}{STATUS.md}.
\end{thebibliography}
\end{document}
